% ===========================================================================
\part{The front end}
% ===========================================================================

\chapter{Pages, routes, and layout}
\label{ch:pages}

\section{Decide the URLs before you write the components}

URLs are the part of a website that is hardest to change later, because other people's links point
at them. They are also the cheapest thing to get right today, since there are none yet apart from
\code{/}.

\begin{table}[htbp]
\centering
\caption{A route plan for this site. The \code{:slug} routes are the reason
Chapter~\ref{ch:datamodel} argues for a \code{slug} column.}
\label{tab:routes}
\small
\begin{tabular}{@{}L{0.24\textwidth}L{0.30\textwidth}L{0.40\textwidth}@{}}
\toprule
\textbf{Path} & \textbf{Reads from} & \textbf{Note} \\
\midrule
\code{/} & \code{people}, \code{contact\_links} & who you are; the only page that must load fast \\
\code{/portfolio} & \code{portfolio\_items} + tags & the list; filterable by tag \\
\code{/portfolio/:slug} & one \code{portfolio\_item} & the detail page; must 404 properly \\
\code{/publications} & \code{publication\_items} & grouped by year, newest first \\
\code{/experience} & \code{job\_experiences} + institutions & the CV, reverse chronological \\
\code{/tags/:slug} & join table & optional; only if tags earn a page of their own \\
\code{*any} & --- & 404, with a real 404 status \\
\bottomrule
\end{tabular}
\end{table}

\begin{pattern}[Slugs, not integers]
\code{/portfolio/egonik-site} is a URL you can put in a CV. \code{/portfolio/7} is a URL that
changes meaning if you ever re-seed the database. Add the \code{slug} column in the same migration
that fixes defect \dnum{1}, before there is any content to migrate.
\end{pattern}

\section{Layout: one shell, many pages}

Every page shares a header, navigation and footer. Leptos expresses that with a nested route: a
parent route renders the shell and an \code{<Outlet/>}, and the children render into it.

\begin{diagram}{The route tree. \code{ParentRoute} renders the shell once; only the
\code{<Outlet/>} contents change on navigation, so the header is never re-rendered and never
flashes.}{fig:routes}
\begin{tikzpicture}[node distance=4mm]
  \node[shared,minimum width=54mm] (shell) {\code{<ParentRoute path="" view=Shell/>}\\[1pt]\scriptsize
    header $+$ nav $+$ \code{<Outlet/>} $+$ footer};
  \node[cli,minimum width=40mm,below left=11mm and 22mm of shell] (home) {\code{path=""}\\\scriptsize \code{<HomePage/>}};
  \node[cli,minimum width=40mm,below=11mm of shell] (port) {\code{path="portfolio"}\\\scriptsize \code{<PortfolioList/>}};
  \node[cli,minimum width=40mm,below right=11mm and 22mm of shell] (pubs) {\code{path="publications"}\\\scriptsize \code{<Publications/>}};
  \node[cli,minimum width=40mm,below=8mm of port] (det) {\code{path="portfolio/:slug"}\\\scriptsize \code{<PortfolioDetail/>}};
  \node[box,fill=wash,minimum width=40mm,below=8mm of pubs] (nf) {\code{WildcardSegment}\\\scriptsize \code{<NotFound/>} $+$ status 404};

  \draw[flow] (shell) -- (home);
  \draw[flow] (shell) -- (port);
  \draw[flow] (shell) -- (pubs);
  \draw[flow] (shell.south) ++(0,-2mm) -- (det.north);
  \draw[flow] (shell) -- (nf);
\end{tikzpicture}
\end{diagram}

\begin{lstlisting}[style=rs,caption={\file{src/app.rs} --- the whole file, and it should stay roughly this size.}]
use leptos::prelude::*;
use leptos_meta::{provide_meta_context, Meta, Stylesheet, Title};
use leptos_router::components::{Outlet, ParentRoute, Route, Router, Routes};
use leptos_router::{ParamSegment, StaticSegment, WildcardSegment};

use crate::ui::layout::Shell;
use crate::ui::pages::{HomePage, NotFound, PortfolioDetail, PortfolioList, Publications, Experience};

#[component]
pub fn App() -> impl IntoView {
    provide_meta_context();

    view! {
        <Stylesheet id="leptos" href="/pkg/egonik-site.css"/>
        <Title text="Eduardo Gonik"/>
        <Meta name="description" content="Personal site: portfolio, publications, experience."/>

        <Router>
            <Routes fallback=NotFound>
                <ParentRoute path=StaticSegment("") view=Shell>
                    <Route path=StaticSegment("")             view=HomePage/>
                    <Route path=StaticSegment("portfolio")    view=PortfolioList/>
                    <Route path=(StaticSegment("portfolio"), ParamSegment("slug"))
                                                              view=PortfolioDetail/>
                    <Route path=StaticSegment("publications") view=Publications/>
                    <Route path=StaticSegment("experience")   view=Experience/>
                    <Route path=WildcardSegment("any")        view=NotFound/>
                </ParentRoute>
            </Routes>
        </Router>
    }
}
\end{lstlisting}

\begin{pitfall}[Use \code{<A>} for internal links, \code{<a>} for external ones]
\code{leptos\_router}'s \code{<A href="/portfolio">} intercepts the click and navigates
client-side --- journey 2 from Chapter~\ref{ch:lifecycle}. A plain \code{<a href="/portfolio">}
triggers a full page load, throwing away the wasm bundle you just downloaded and re-running SSR.
Both work; one is much slower. For links off your site, plain \code{<a>} is correct.
\end{pitfall}

\section{Per-page titles, and the SEO detail that is easy to miss}

\code{leptos\_meta}'s \code{<Title>} works inside any component, including a page, and the last one
rendered wins. So each page sets its own:

\begin{lstlisting}[style=rs]
#[component]
pub fn Publications() -> impl IntoView {
    view! {
        <Title text="Publications | Eduardo Gonik"/>
        // ...
    }
}
\end{lstlisting}

There is a subtlety when the title depends on \emph{loaded data} --- a portfolio detail page whose
title is the project name. During server rendering, the HTML \code{<head>} is sent before the body
finishes. If the title comes from a normal \code{Resource}, the head may already be on the wire by
the time the data arrives, and the crawler sees the generic title.

\begin{pattern}[\code{Resource::new\_blocking} for data that feeds the head]
\code{Resource::new\_blocking} has the same signature as \code{Resource::new} but holds the HTTP
response until the data has loaded. Use it for exactly the one resource whose value determines
\code{<Title>}, \code{<Meta>} or the HTTP status --- and for nothing else, because it costs you
streaming.
\begin{lstlisting}[style=rs]
// the project name is in the title and the meta description -> blocking
let item = Resource::new_blocking(move || slug(), |s| get_portfolio_item(s));
\end{lstlisting}
\end{pattern}

\section{404 must be a real 404}

The starter template already gets this right --- verifiably, and this is worth checking rather than
assuming, because the mechanism is easy to break later.

\begin{verified}[Verified against the running binary]
\begin{lstlisting}[style=out]
$ curl -s -o /dev/null -w '%{http_code}\n' localhost:3000/nonexistent-path
404
$ curl -s localhost:3000/nonexistent-path | grep -o '<h1>[^<]*</h1>'
<h1>Not Found</h1>
\end{lstlisting}
\end{verified}

So the 404 path works today. One subtlety, though: \file{app.rs} declares
\emph{both} a wildcard route and a \code{Routes} fallback:

\begin{lstlisting}[style=rs]
<Routes fallback=move || "Not found.">                    // <-- unreachable
    <Route path=StaticSegment("") view=HomePage/>
    <Route path=WildcardSegment("any") view=NotFound/>    // <-- matches everything else first
</Routes>
\end{lstlisting}

The wildcard route matches every unmatched path, so the fallback string is dead code that will
never render. That is harmless and misleading in equal measure --- it looks like the 404 handler,
and it is not. Delete it, or point it at \code{NotFound} so that the two agree.

\code{NotFound} sets the status inside a \code{cfg} block:

\begin{lstlisting}[style=rs]
#[cfg(feature = "ssr")]
{
    let resp = expect_context::<leptos_actix::ResponseOptions>();
    resp.set_status(actix_web::http::StatusCode::NOT_FOUND);
}
\end{lstlisting}

The \code{cfg} is not defensive noise --- \code{ResponseOptions} exists only during server
rendering, because only then is there an HTTP response to modify. If a visitor navigates to a bad
URL \emph{after} hydration there is no new request to set a status on, which is a real and
unavoidable asymmetry: the first load of a missing page returns 404, a client-side navigation to one
does not. That is correct behaviour and matches how every SPA works.

The same mechanism applies to a \code{/portfolio/:slug} that does not exist, and it is the reason
Chapter~\ref{ch:errors} insists on calling \code{set\_status} in the error mapping layer. A detail
page for a deleted project must return 404, or search engines will index it forever.

% ===========================================================================
\chapter{Loading data into a page}
\label{ch:loading}

\section{The three resource types, and which one you want}

This is the most common place to make a wrong choice that appears to work. All three types load
async data; they differ in whether they participate in server rendering.

\begin{table}[htbp]
\centering
\caption{Resource types in Leptos 0.8.}
\label{tab:resources}
\small
\begin{tabular}{@{}L{0.19\textwidth}cL{0.30\textwidth}L{0.34\textwidth}@{}}
\toprule
\textbf{Type} & \textbf{Runs in SSR?} & \textbf{Constructor} & \textbf{Use it when} \\
\midrule
\code{Resource} & \textcolor{moss}{yes} & \code{new(source, fetcher)} & the default. Reactive to
its source; result is serialised into the HTML \\
\code{Resource\allowbreak::new\_blocking} & \textcolor{moss}{yes} & same & the value determines
\code{<Title>}, \code{<Meta>} or the HTTP status \\
\code{OnceResource} & \textcolor{moss}{yes} & \code{new(future)} & loads exactly once, never
re-fetches, no source signal \\
\code{LocalResource} & \textcolor{crimson}{\textbf{no}} & \code{new(fetcher)} & browser-only APIs,
or a future that is not \code{Send} \\
\bottomrule
\end{tabular}
\end{table}

\begin{pitfall}[\code{LocalResource} is the trap]
Its documentation says it ``will only begin loading data if you are on the client''. Reach for it
--- perhaps because a \code{Send} bound was inconvenient --- and your page silently stops being
server-rendered: the HTML ships as an empty shell, the content appears only after the wasm loads,
search engines see nothing, and none of it produces a warning. If you find yourself reaching for
\code{LocalResource} to make a \code{Send} error go away, the real problem is almost always a
non-\code{Send} value held across an \code{.await} inside the server function --- typically a
database connection. Fix that instead; Chapter~\ref{ch:blocking} shows how.
\end{pitfall}

\begin{pitfall}[A \code{Send} error that points at the wrong line]
\code{Resource::new} requires the fetcher's future to be \code{Send} and \code{T} to be
\code{Send + Sync + Serialize + DeserializeOwned}. If your server function holds a
\code{PooledConnection} across an \code{.await}, the resulting \code{Send} violation surfaces at
\code{Resource::new} --- in a different file from the actual mistake. When a \code{Send} error names
a resource, look inside the server function it calls.
\end{pitfall}

\section{The canonical read pattern}

\begin{pattern}[\code{Resource} $+$ \code{Transition} $+$ \code{ErrorBoundary} $+$ \code{Suspend}]
\begin{lstlisting}[style=rs,caption={\file{src/publications/components.rs} --- the shape to grow the
current placeholder into.}]
use leptos::prelude::*;
use leptos::either::Either;
use crate::publications::api::list_publications;

#[component]
pub fn Publications() -> impl IntoView {          // NOTE: pub
    let publications = Resource::new(|| (), |_| list_publications());

    view! {
        <Transition fallback=move || view! { <p class="text-slate-500">"Loading..."</p> }>
            <ErrorBoundary fallback=|errors| view! {
                <p class="text-red-700">
                    "Could not load publications: "
                    {move || errors.get()
                        .into_iter()
                        .map(|(_, e)| e.to_string())
                        .collect::<Vec<_>>()
                        .join("; ")}
                </p>
            }>
                {move || Suspend::new(async move {
                    publications.await.map(|rows| if rows.is_empty() {
                        Either::Left(view! { <p>"Nothing published yet."</p> })
                    } else {
                        Either::Right(rows.into_iter()
                            .map(|p| view! { <PublicationCard publication=p/> })
                            .collect_view())
                    })
                })}
            </ErrorBoundary>
        </Transition>
    }
}
\end{lstlisting}
\end{pattern}

Five details in that snippet are load-bearing:

\begin{description}
  \item[\code{Transition}, not \code{Suspense}.] \code{<Suspense>} shows its fallback every time
    any nested resource reloads, so a re-fetch makes the list flash back to ``Loading\ldots''.
    \code{<Transition>} shows the fallback on first load only and keeps the stale content visible
    while new data arrives. For lists that refetch, \code{Transition} is almost always what you
    want.

  \item[\code{ErrorBoundary} goes \emph{inside} the transition.] Not outside. The upstream Leptos
    example still carries a comment noting that the outside-in arrangement breaks \code{Suspense}
    under Actix. Inside-out works; take the ordering as given.

  \item[\code{Suspend::new} lets you \code{.await} the resource in the view.] The alternative is
    \code{publications.get()} returning \code{Option<Result<...>>} and a nest of matches.
    \code{Suspend} removes the \code{Option}, because inside a suspense boundary the framework
    guarantees the value is there.

  \item[\code{Either::Left} / \code{Either::Right}.] Two \code{view!} invocations produce two
    \emph{different static types} in Leptos 0.8, so an \code{if}/\code{else} that returns markup
    will not compile without wrapping the arms. \code{Either} is the cheap fix;
    \code{.into\_any()} is the type-erasing one, slightly slower and sometimes clearer.

  \item[Errors surface by themselves.] A \code{Result} rendered in a view renders the \code{Ok}
    value and, on \code{Err}, renders nothing while registering the error with the nearest
    \code{ErrorBoundary}. You do not match on the error --- you let it propagate to the boundary.
\end{description}

\section{Writes: \code{Action}, not \code{Resource}}

You have no writes yet. When you add them, the pairing is strict.

\begin{antipattern}[A mutation inside a \code{Resource}]
A \code{Resource} runs during server rendering \emph{and} may run again on the client. Put an
\code{INSERT} in one and you will insert twice, and the second row will appear only in production
where the timing differs. Reads are Resources; writes are Actions.
\end{antipattern}

\begin{lstlisting}[style=rs,caption={A write, and the read that refreshes itself when the write lands.}]
let add_item = ServerAction::<AddPortfolioItem>::new();

// the resource re-runs whenever the action's version changes
let items = Resource::new(
    move || add_item.version().get(),
    |_| list_portfolio(),
);

view! {
    <ActionForm action=add_item>
        <input type="text" name="title"/>
        <button type="submit">"Add"</button>
    </ActionForm>
}
\end{lstlisting}

\begin{why}[Why \code{<ActionForm>} rather than an \code{on:click} handler]
\code{<ActionForm>} degrades to an ordinary HTML \code{<form method="POST">}. With no JavaScript it
still submits, the server function still runs, and \code{leptos\_actix} notices that the request
accepts \code{text/html} and responds with a 302 back to the referring page. You get progressive
enhancement free. A click handler that calls the server function directly does not.
\end{why}

\begin{pitfall}[\code{Vec} parameters need \code{\#[server(default)]}]
The default URL encoding omits an empty \code{Vec} entirely, and deserialisation on the server then
fails with a missing-argument error at runtime --- not at compile time. Any server function taking
a \code{Vec} needs the annotation:
\begin{lstlisting}[style=rs]
#[server]
pub async fn set_tags(#[server(default)] tags: Vec<String>) -> Result<(), ServerFnError> { ... }
\end{lstlisting}
\end{pitfall}

\section{One thing to internalise about server functions}

\begin{keyidea}
A server function is a public HTTP endpoint. Anyone can \code{curl} it. It is not a private
mechanism for your front end to talk to your back end.
\end{keyidea}

The Leptos book puts it as ``server functions are not magic; they're syntax sugar for defining a
public API''. Two practical consequences for this site:

\begin{itemize}
  \item \code{portfolio\_items} has a \code{public} boolean. A server function called
    \code{list\_portfolio} must filter on it \emph{in the query}, not in the component. Filtering in
    the component leaves the private rows in the JSON payload, visible in the network tab and in the
    server-rendered HTML.
  \item Every write function needs its own authorisation check, in its own body. There is no
    middleware that will do it for you by virtue of the function being ``internal'', because it is
    not internal.
\end{itemize}

% ===========================================================================
\chapter{Tailwind CSS}
\label{ch:tailwind}

\section{What you are switching from and to}

\file{style/main.scss} is three lines of Sass. cargo-leptos compiles it with a \code{sass} binary it
downloaded itself, and injects the result at \file{/pkg/egonik-site.css} --- which is already what
\file{src/app.rs} references, so that part needs no change.

The good news is that the Tailwind side needs no Node, no \file{package.json}, and no \code{npm
install}. cargo-leptos manages the Tailwind CLI exactly as it manages Sass.

\begin{verified}[Verified: the Tailwind CLI is already on this machine]
\begin{lstlisting}[style=out]
$ ls ~/.cache/cargo-leptos/
  sass-1.86.0/  tailwindcss-v4.2.1/  wasm-bindgen-0.2.125/  wasm-bindgen-0.2.126/
\end{lstlisting}
cargo-leptos 0.3.6 hardcodes \code{v4.2.1} as its default Tailwind version --- so this is
\textbf{Tailwind 4}, not 3, and the difference matters for every instruction below. The binary is
downloaded from the \code{tailwindlabs/tailwindcss} GitHub releases; \code{LEPTOS\_TAILWIND\_VERSION}
overrides the version if you ever need to.
\end{verified}

\section{The change, in four steps}

\begin{runbook}[Adopting Tailwind 4]
\textbf{1. Create the input file.} One line is genuinely sufficient in Tailwind 4:
\begin{lstlisting}[style=sh]
# style/tailwind.css
@import "tailwindcss";
\end{lstlisting}

\textbf{2. Point cargo-leptos at it}, in \file{Cargo.toml}:
\begin{lstlisting}[style=toml]
[package.metadata.leptos]
tailwind-input-file = "style/tailwind.css"
# style-file = "style/main.scss"     <-- delete this line (see below)
# do NOT add tailwind-config-file    <-- see the pitfall below
\end{lstlisting}

\textbf{3. Delete \file{style/main.scss}} and the three lines in it. Tailwind's preflight already
sets sensible defaults, and the \code{text-align: center} in that file is not something you want
site-wide.

\textbf{4. Rebuild.} \code{cargo leptos watch}. The href in \file{app.rs} does not change: the
combined output is always \file{<site-root>/<site-pkg-dir>/<output-name>.css}, which for this
project is \file{target/site/pkg/egonik-site.css}, served at \code{/pkg/egonik-site.css}.
\end{runbook}

\section{The traps, in order of how much time they cost}

\begin{pitfall}[Do not set \code{tailwind-config-file}]
This is the highest-value warning in the chapter. Under cargo-leptos 0.3.6 the default Tailwind is
v4, and \textbf{the v4 CLI has no \code{--config} flag at all}. It accepts the argument silently,
exits 0, and ignores it --- even if the path does not exist.

Worse: if you set the key and the file is missing, cargo-leptos \emph{writes a v3-style
\file{tailwind.config.js} into your project for you}, complete with \code{content} globs. That file
then sits in your repository looking authoritative while having no effect whatsoever. When purging
misbehaves, you will edit it, and nothing will happen.

Setting \code{tailwind-config-file} without \code{tailwind-input-file} is also a hard error, with
the message \code{The Cargo.toml `tailwind-input-file` is required when using `tailwind-config-file`]}
--- the stray bracket is really in the source.
\end{pitfall}

\begin{pitfall}[A \code{tailwindcss} on your PATH silently wins]
cargo-leptos calls \code{which::which("tailwindcss")} first and uses whatever it finds,
\emph{with no version check}. If a Tailwind 3 CLI is installed globally, you build with v3 while
cargo-leptos still takes its v4 code path and omits \code{--config} --- so you get v3 with no
configuration, which purges nearly everything. The same happens if a \file{package.json} exists at
the workspace root, because cargo-leptos prepends \file{node\_modules/.bin} to \code{PATH}.

\file{end2end/package.json} exists in this repository. It is in \file{end2end/}, not the workspace
root, so it does not currently trigger this --- but if you ever add a root-level
\file{package.json}, remember this paragraph. Diagnose with \code{cargo leptos watch -vv} and read
the logged \code{tailwindcss --input ...} line.
\end{pitfall}

\begin{pitfall}[\file{target/} must stay in \file{.gitignore}]
Tailwind 4 detects classes by scanning the filesystem, and it skips whatever \file{.gitignore}
excludes. \file{target/} contains the previous CSS output and vendored crate sources; scanning it is
slow --- one reported case went from 2\,ms to 54\,s --- and it can resurrect classes you deleted,
making purging look broken. Your \file{.gitignore} lists \file{target/} already. Leave it there.
\end{pitfall}

\begin{pitfall}[Three different path bases in one pipeline]
This is the usual cause of ``my classes are being purged''.
\begin{itemize}
  \item \code{tailwind-input-file} resolves against the directory of the \file{Cargo.toml} holding
    the metadata table.
  \item \code{@source} directives inside the CSS resolve against \emph{the CSS file}.
  \item Tailwind's automatic detection resolves against \emph{the process working directory}, and
    cargo-leptos \code{chdir}s to the cargo workspace root.
\end{itemize}
For this single-crate project all three coincide, so it works with no configuration. If you ever
move to a workspace, pin it explicitly and stop relying on the coincidence:
\begin{lstlisting}[style=sh]
@import "tailwindcss" source("../src");
\end{lstlisting}
\end{pitfall}

\section{Class detection inside \code{view!} macros}

The question everyone asks: does Tailwind find classes written inside a Rust macro? Yes, without
configuration.

\begin{verified}[Verified against the cached v4.2.1 binary]
With a bare \code{@import "tailwindcss";} and no \code{@source}, all of these forms produced their
rules in the output:
\begin{lstlisting}[style=rs]
view! {
    <div class="flex gap-4 text-sky-500">
        <p class:underline=true class:font-bold=cond>"hi"</p>
        <span class=("bg-red-500", move || active())>"t"</span>
    </div>
}
\end{lstlisting}
Note in particular that \code{class:underline=true} yields \code{underline}. The v3-era workaround
--- a JavaScript \code{transform} that stripped \code{class:} prefixes --- is unnecessary under v4,
and there is no \code{transform} option in v4 anyway.
\end{verified}

\begin{antipattern}[Computing class names at runtime]
\begin{lstlisting}[style=rs]
let colour = if item.public { "green" } else { "red" };
view! { <span class=format!("text-{colour}-600")>...</span> }   // BROKEN
\end{lstlisting}
Tailwind scans \emph{source text}. \code{text-green-600} never appears in your source, so the rule
is never generated and the element is unstyled. Write both class names out in full and choose
between them:
\begin{lstlisting}[style=rs]
let colour = if item.public { "text-green-600" } else { "text-red-600" };
view! { <span class=colour>...</span> }                        // works
\end{lstlisting}
This is the single most common Tailwind bug in any language, and it is easier to hit in Rust because
string formatting is so natural.
\end{antipattern}

\section{Two subtleties worth knowing before they surprise you}

\subsection{If you keep the Sass file, it will override Tailwind}

\code{style-file} and \code{tailwind-input-file} are \emph{not} mutually exclusive. cargo-leptos
runs Sass and Tailwind as two independent processes and simply concatenates the results, Sass first,
then Tailwind.

That ordering suggests Tailwind wins ties. It does not. Tailwind 4 wraps everything in real cascade
layers (\code{@layer theme, base, components, utilities;}), and \textbf{unlayered CSS beats layered
CSS} in the cascade regardless of order. So your hand-written Sass will override Tailwind utilities
of equal specificity even though it appears first in the file. Two consequences:

\begin{itemize}
  \item If you keep both, wrap your own CSS in a layer so it participates in the cascade normally.
  \item \code{@apply}, \code{@theme} and \code{@plugin} work \emph{only} in the Tailwind input file.
    Put \code{@apply} in the \file{.scss} and dart-sass passes it through untouched --- you get
    invalid CSS in the output rather than an error.
\end{itemize}

The simplest resolution is the one in the runbook: drop \code{style-file} entirely and put
everything, including \code{@theme} customisation, in \file{style/tailwind.css}.

\subsection{\code{browserquery = "defaults"} downlevels modern CSS}

cargo-leptos runs the combined CSS through \code{lightningcss} on \emph{every} build --- not only
release --- using \code{browserquery} as its target list. Tailwind 4 output depends on
\code{@property}, \code{oklch()} and \code{color-mix()}. With a broad target list like the current
\code{"defaults"}, lightningcss will rewrite that output for old browsers, and colours or custom
properties can come out subtly wrong.

\begin{pattern}[Narrow \code{browserquery}]
\begin{lstlisting}[style=toml]
browserquery = "last 2 versions and not dead"
\end{lstlisting}
There is no way to disable the lightningcss pass. Note also that if lightningcss cannot parse
Tailwind's output, cargo-leptos 0.3.6 falls back to emitting the unminified CSS with only a
\emph{trace}-level log --- so a real problem is invisible unless you run with \code{-vv}. Older
versions failed the build outright, which was arguably friendlier.
\end{pattern}

\section{Watching and reloading}

\begin{itemize}
  \item \textbf{Editing a \code{.rs} file does re-run Tailwind.} cargo-leptos knows that CSS can
    come from source files when \code{tailwind-input-file} is set, so a Rust change triggers a
    style rebuild. Adding a utility class to a \code{view!} works in \code{watch} mode.
  \item \textbf{A style-only change reloads just the CSS}, over the live-reload websocket, without
    a full page refresh.
  \item \textbf{But \code{--hot-reload} cannot add new classes.} The view-patching hot reload
    updates markup in place without a rebuild, and a brand-new Tailwind class needs a rebuild to
    exist. Expect to restart occasionally.
  \item \textbf{Only the exact input-file path is watched}, not files it \code{@import}s. If you
    split your CSS across several files, list the others under \code{watch-additional-files}.
    There is also a known open bug where a change to the input file is detected but the new CSS is
    not applied until a manual restart.
\end{itemize}

\section{Design tokens, once Tailwind is in}

A brief note, because it is the thing that makes a Tailwind codebase pleasant or awful. In v4,
customisation happens in CSS rather than JavaScript:

\begin{lstlisting}[style=sh,caption={\file{style/tailwind.css} --- tokens defined once, used everywhere.}]
@import "tailwindcss";

@theme {
  --color-ink:   oklch(0.22 0.02 250);
  --color-paper: oklch(0.98 0.005 90);
  --color-accent: oklch(0.55 0.15 40);
  --font-display: "Libertinus Serif", Georgia, serif;
}
\end{lstlisting}

Those become utilities automatically: \code{text-ink}, \code{bg-paper}, \code{border-accent},
\code{font-display}. Define the palette once, then never write a raw hex value in a \code{view!}
again --- that discipline is what keeps a hand-built site from drifting into fifteen shades of
almost-grey.
